Materials innovation is one of the main driving force behind the modern socioeconomic advances. It encompasses multiple disciplinaries that impact directly upon our daily life. For examples, the development of microelectronics with low-power consumption and fast processing speeds,\cite{lin2019two} electrocaloric materials\cite{valant2012electrocaloric} that enables the effective harvesting of waste heats in microelectronics, next-generation materials that can effectively convert renewable energy sources into electricity,\cite{zhang2022renewed,han2022ferroelectric} catalysts for the large-scale synthesis of industrially relevant chemicals,\cite{chang2021emerging} new drug molecules for the therapeutic treatments of existing and emerging life-threatening diseases,\cite{tanoli2025computational} as well as materials that could sustain and perform under extreme environments (such as inside nuclear reactors or at the outer space).\cite{mondal2021thermal,tadjer2022toward} There has been a long quest to achieve the goal of \emph{dial-a-material} to combat the increasing socioeconomic demands, where one could design the chemistry and structure, followed by the successful synthesis, of new materials with pre-defined chemical/physical properties for specific application, but the  destination for such a quest is still seems to on a far reaching horizon of an uncharted land. The pursuit for a room-temperature superconductor for over a century,\cite{pickett2023colloquium} is probably one of the most compelling and representative examples.  

The challenges behind materials discoveries are multilevelled. 
At the first level, there is a combinatorially large space to explore, in which different chemical elements and the atomic configurations can be combined together to give new materials,\cite{reymond2015chemical,groom2016cambridge,draxl2019nomad} even one builds these structures by following some very simple chemical and structural rules, such as the octet rules and known crystal symmetries.
At the second level, even with a given chemical composition, material can exhibit multiple different structures (polymorphs) with widely varying thermodynamic stabilities.\cite{doye1999evolution,engel2018mapping,senocrate2019thermochemical,yu2010polymorphism,bauer2001ritonavir} To include the next level of complexity, the thermodynamic  stability for materials itself is a complex, multivariable metric to determine both experimentally and theoretically, which is tightly coupled with the occurrence of defects and disorders that heavily depends on the synthesis condition. Whilst basic chemical intuitions based on elementary bonding theory is often very useful for us to rationalise the stability of a material, compelling evidences do suggest that more exotic quantum phenomena can also play an important role in materials' structural stabilities, as in the case of ferroelectric (in)stability in perovskites.\cite{bersuker2013pseudo} This naturally leads to the final, fourth level of challenge in our ranking, in which many material properties are the manifests of the responses from complex many-body quantum interactions in the material as a response to the external stimuli.\cite{slack1973nonmetallic,allen1993thermal,Migdal1958,Eliashberg1960,giustino2017electron,donchev2006many,ramasubramaniam2012large} It means that many of the exotic phenomena that drives the current wave of materials innovations may not be understood with simple physical laws and equations such as the classical Newtonian mechanics or Maxwell's theory of electromagnetism.
Certainly, one may rank these challenges in a different order, depending on the perspective one tends to approach this problem.

Whilst our technological advancements are largely driven by experimental breakthroughs in materials discoveries of new materials and phenomena that are not existing in the current material space, the interwinding complexities discussed above are often bringing us tremendous challenges to understand the experimental results, derive useful unbiased generalisations, from which predictive decisions for materials discoveries can be made. As such, computational materials, particularly atomistic modellings, have been playing a pivotal rule in dissecting these complex factor, to deepen our understanding behind the fundamental materials physics and chemistry from the atomistic level. The main benefit of computational modelling, is to provide an pragmatic approach, which we can draw analogy to Taylor expansions in mathematics or perturbation theory in quantum mechanics, that enables us to isolate and examine how different (sets of) physical effect(s) contribute, either individually or corporately, to the new behaviours observed in experiments. Such a fine-grained control, that separate and isolate different physical contributions towards a specific observed phenomena may not always be achievable experimentally. 

At the most fundamental level, quantum mechanics outline the principles that govern the interactions among the nuclei and electrons, that lead to the  structures, stabilities and the properties of materials at the atomic scale. Nevertheless, the many-body quantum problems\cite{mahan2000many} are inherently complex, and often impossible to be solved, to truly reflect the behaviours of real materials in both the time and length scales that we are commonly dealing with. On the other hand, to link microscopic phenomena with macroscopic observables, we also need to apply tools from statistical mechanics to compute the thermodynamically relevant behaviours of materials. As a small divergence at this point, although the later merely compounded the challenges that we are facing, we must point out that both theoretical frameworks have a strong linkage with probabilistic theory, in which \emph{sampling} played a critical role in helping to establish the computational approaches to tackle problems in their respective fields, and tight strongly with modern machine learning (ML) and artificial intelligence (AI),\cite{murphy2012machine} which we shall elaborate further in this monograph.

To overcome the many-body challenges, theoreticians have long been working on the developments of plausible \emph{approximations} \cite{harrison2018review,kohn1999nobel,golze2019gw} to solve these problems, to which, now materials modelling practitioners have the access to a rich repository of practically feasible simulation methodologies to compute the structures, energetics and the associated properties of materials with predictive accuracies, enabling many new materials to be predicted before more expensive laboratory synthesis and characterisation efforts are devoted to minimise the costs in materials discoveries.\cite{Lv2015Weyl,pulido2017functional,Mounet2018Exfoliation,Sun2019TernaryNitrides} Some of these methods, which has profoundly shaped the current landscape of the materials modelling practice, and shall be discussed in more details in this monograph, include: (1) Empirical atom-atom potentials or forcefields,\cite{harrison2018review} that offered simple parameterised analytical expressions to compute the interaction energies for atomistic systems, the many fundamental concepts in which underpinned the development of mathematical representations of atomistic structures for modern ML.\cite{Musil2021StructuralRepresentations} (2) Density Functional Theory (DFT),\cite{kohn1999nobel} which unarguably one of the most successful approximations developed for treating the many-body electronic structure problem that is the main workhorse for computational chemistry and materials science. (c) Many-body methods, such as GW approximations\cite{golze2019gw} for materials, which provided us to compute materials' properties at even better accuracies. These are ultimately integrated with (d) statistical sampling methods that enable us to   sample  the potential energy surface (PES) to predict materials' behaviour that are thermodynamically relevant to experimental observations. On the other hand, the availability of these methods, require materials modelling practitioners\footnote{which, for this monograph, the name refer to those who are not directly involved in the development of new theories or methods, but primarily focusing on the applications of the established methods in materials simulations}  to have a good solid understanding behind the fundamental principles and assumptions adopted behind these methodologies, as well as the trade-off between their accuracies, computational costs and system sizes capable of dealing with, in order to appropriately apply them in real problem solving.

From the physics point of view, the uniting concept behind all these developed methodologies is the energy of the atomistic system, that enters either as the  output of the model (such as empirical potentials), one of the key quantity that many others (electron densities, response properties) are associated with or can be derived from (such as in DFT), or the primary quantity to be sampled in statistical samplings. 

Therefore, from both the realistic materials application and theoretical method development perspectives, surveying the PES (approximated at a chosen theory level) of a given material system, as a function of its atomistic structures $E(\bm{R})$, under some prescribed thermodynamic conditions, can be considered as a central theme of computational materials science. The determinations of other associated physical properties, such as the anharmonic lattice dynamics and thermal conductivities,\cite{knoop2024tdep,tadano2015self,phono3py} dielectric responses,\cite{souza2002first} magnetic ground states,\cite{fender2025altermagnetism} as well as electron-phonon interactions,\cite{giustino2017electron} can all be cast into some form of sampling problems in more complex manifolds, such as on a PES in $\mathbb{R}^{3N}$ (with $N$ being the number of atoms) that are approximated beyond a harmonic truncation of the Taylor expansion of $E(\bm{R})$ in real space or alternatively as the combinatorial couplings among different harmonic phonons in the reciprocal space, the inclusion of the spin degree-of-freedoms in the PES, or the coupling between electron and phonon degrees of freedoms in which the PES for both the electronic and phononic subsystems need to be sampled jointly. A more concrete example lies in the study of heterogeneous electrocatalysts in explicit solvents,\cite{Andrade2020ProtonTransfer,Qin2023CO2Electroreduction,Sakong2018ElectricDoubleLayer} which has become increasingly popular with the advance in computational power, that one finds many of the aforementioned problems meet in a single system at the intersection between chemistry and physics.

The increase in the computational power, that is backed by the developments of high performance computers (HPC), and more recently, Graphical Processing Units (GPU), not only enriched the individual's capabilities in tackling computational materials problems with greater complexities, but also the communal developments of automated workflow management research software\cite{ganose2025atomate2,ong2013python} to facilitate the high-throughput, multi-staged computational modelling over families of materials systems. This ultimately leads to the establishment and continuous growths of several community-driven large-scale computational materials databases\cite{jain_apl_mater,draxl2019nomad,kirklin2015open,tran2023open} over the past decade, that not only has become useful centralised repositories for exploring materials' stabilities, structures and various physical properties (such as electronic structures, phonon properties, and magnetic ground states), but most importantly, provided the much needed community standard for materials modelling practices to enhance the consistency, reproducibility and credibility of the new research outputs in similar area.

It is probably not a surprise that this community achievement comes at the same time when ML and AI become increasingly popular, and start to reshape the landscape of computational research in chemistry, physics and materials science. Fundamentally, ML and AI are data-driven and statistical approaches, that enable us to extract complex hidden patterns (knowledges) from the dataset,\cite{glielmo2021unsupervised} that may be further harnessed to build various different models\cite{deringer2021gaussian,behler2021four,duval2023hitchhiker} to accelerate the computation and predictions of materials' stabilities and properties. Lately,  the use of generative AI\cite{zeni2025generative} has opened up the real hope for us to predict new chemistries and structures for materials that are not yet known, which it shall be demonstrated in this monograph, that this new frontier has offered us a new paradigm to reformulate transitional statistical samplings methods, making them to become better guided by the (equilibrium) probabilistic distributions of materials.

Whilst the developments in ML and  AI have brought many excitements into the field, a key challenge behind the adaptation of these breakthroughs in computer science into molecular and materials domain, is to appropriately incorporate the fundamental principles in atomistic modellings as discussed above, perhaps most critically, \emph{symmetries}, into these models mathematically.\cite{bartok2013representing,musil2021physics,grisafi2018symmetry,duval2023hitchhiker,nigam2022unified} This has shown to be important for enhancing the models' data efficiency and extrapolability, regardless whether a simple ML surrogate or more complex AI model is to be used. There has been a long criticism that DFT has often been used blindly as `black-box' methods by the practitioners,\cite{kruse2012standard} and the shear mathematical complexities in the newly developed AI models across the community can impose similar risks, with models being chosen by the practitioners based on popularity (or more bluntly put, number of citations), without properly understand the `nitty-gritty' behind the models, especially how and why two models differ from each other.

To help overcoming this `black-box' perceptions, and facilitate the future computational materials modelling practitioners to understand the multiple connected layers in the physical models, their workable approximations, probability theory in statistical thermodynamics and modern AI model developments, as well as the vibrant field of group and representation theory that played a pivotal role in atomistic ML and AI, I believe it will be valuable in writing a monograph that is based on a centralised concept of statistical sampling, to provide the readers with an overview on the topics discussed above in a pedagogically approachable way. This contents and organisations behind this monograph is building on my own research and teaching experiences. In particular, key mathematical equations will be presented to retain the necessary theoretical rigorous, but I also tend to provide a more intuitive, hitch-hiker style interpretations behind physical meanings behind these equations. I found that such an approach will be particularly useful in understanding the latest generative AI models for molecules and materials, which have extensively borrowed  mathematical tools from differential geometries, group theory and stochastic processes, which some readers may found intimidated to comprehend.

Hence, the primary objectives behind the present monograph, is to facilitate the practitioners in atomistic materials modelling, to: Firstly, establish the necessary physical and mathematical foundations for atomistic modelling. Secondly, establish an appreciation behind the energy-structure-property relationships for materials, that underpinned the developments of physics-inspired (structural) material representations, that significantly shaped the current landscape in ML and AI-driven materials science.  Thirdly, establish the linkage among traditionally applied simulation method and modern data-driven approaches, with the former acting more than simply a data-generator for the latter, but more fundamentally and conceptually linked through statistical thermodynamics and stochastic sampling, which we hope can facilitate the reader to establish a more unified viewpoint in understand the methodologies that have been developing in the field. Ultimately, these knowledge will be essential in helping the practitioners to better identify the limitations behind the currently available modelling methods, and more importantly, a foundation to advance themselves into the more challenging endeavours of developing  new materials theory and computational methods to tackle emerging problems in materials physics and chemistry, as well as incorporating new AI methodologies in the tackling of these problems. 

Therefore, I hope that both the competent and newly entered practitioners in atomistic modelling can both find inspirations in this, somehow personal, account of the essential knowledges and developments that shaped the modern atomistic materials modelling.

This monograph is organised as following. \cref{chapter:energy_landscape_concept} introduces the concept of energy landscape and surveyed a variety of examples in which the concept is relevant to the problems in materials science. 
\cref{chapter:energy_models} reviews empirical forcefields (or atom-atom potentials), DFT and many-body GW methods, which we treat primarily as the `energy calculators', on the theoretical level of which the PES will be established.  
\cref{chapter:sampling_methods} discussed the traditional statistical sampling methods that includes Monte Carlo, molecular dynamics and other related enhanced sampling methods.  
\cref{chapter:physical_properties} examines approaches with which the structural, dynamic, electronic and magnetic properties can be determined. The modern approaches for studying the electron-phonon interactions is also briefly reviewed to showcase the capabilities for us to tackle many-body physics in materials science that give rise to a variety of exotic physical properties, including superconductivities, and thermoelectric effects. 
\cref{chapter:high-throughput} highlights the developments of high-throughput computational workflows and the large scape computational materials databases. 
\cref{sect:ML_AI_Mat_Sci} develops the foundations and applications of ML and AI models in molecular and materials simulations, in which we pay special focus on how symmetry is incorporated into the model by respecting the invariant and equivariant properties in the atomistic structures through the tools of group and representation theory.
The monograph concludes in \cref{sec:conclusion_outlook}, which some open challenges and future directions will be discussed.
